\chapter{Chapter 4} \label{ch:Chapter4}
\section{Internetworking avanzato}
\subsection{Interdomain Routing}
L'Internet è organizzato in sistemi autonomi (AS), ciascuno dei quali è sotto il controllo di una singola entità amministrativa. Il sistema autonomo (AS) corrisponde a un dominio amministrativo (es.: università, azienda, rete dorsale, ecc.), mentre l’entità amministrativa è colei che assegna le regole politiche e omogenee da far rispettare all’interno della rete affinché essa funzioni correttamente (es.: tutti i router devono essere di un certo tipo, oppure tutti devono usare un certo protocollo come RIP o altri, oppure si misurano i ritardi o i throughput, ecc.). La rete interna di un’azienda potrebbe essere un singolo AS, così come la rete di un singolo fornitore di servizi Internet. Quindi un AS, visto come un insieme di router e reti locali (in figura le zone racchiuse dalle linee colorate), è una sottorete di Internet.
\begin{example}
Le AS possono anche essere i singoli provider di un certo servizio che inglobano le sottoreti dei clienti del servizio offerto (es.: la rete rosa) e che a sua volta sono inglobate e messe in comunicazione con reti nazionali o internazionali più grandi di livelli superiori (es.: la rete verde).
\end{example}
Il GARR (Gruppo per l’ Armonizzazione della Rete della Ricerca) è il consorzio italiano di interconnessione delle reti di ricerca, istruzione e cultura. È una sorta di provider per reti di ricerca ed è parallelo ad altre reti nazionali e internazionali.  Se ci si collega al sito della GARR è possibile vedere in tempo reale le prestazioni e le caratteristiche di tali punti e linee, nonché il carico del traffico (es.: di un certo collegamento si può visualizzare quanti link dedicati ha e con quali performance, se vi sono link di backup e quanto traffico stanno facendo in quel momento). Per quanto riguarda i link Géant, NaMeX, Cogent, ecc. sono i collegamenti verso le reti internazionali o a altre reti nazionali simili a questa italiana. Essi servono a uscire dalle AS e a fare in modo di poter visualizzare ad esempio siti web di altre nazioni o continenti. Affianco la rete europea, di cui il GARR è stilizzato ai soli nodi e linee più rilevanti. Sotto, invece, la rete mondiale. Da notare il link (NaMeX) da Roma verso l’esterno dell’Italia fino in America attraverso l’oceano e che costituisce il punto di accesso/uscita dei pacchetti da fuori/dentro l’Italia. È da ricordare che questa è solo la struttura internet per la reti di ricerca, studio e cultura e che è parallela a altre reti mondiali (es.: quella commerciale).
\subsection{Route Propagation}
Per quanto riguarda il routing, a grandi dimensioni non è possibile utilizzare i protocolli e le tecniche viste finora, in quanto è praticamente impossibile mettere d’accordo tutti i router su una metrica di distanza (costi) e sulle decisioni di instradamento dei pacchetti, che non sono necessariamente sempre guidate da motivi di ottimizzazione dei costi, ma molto spesso sono ragioni “politiche” (= policies, es.: una rete deve essere evitata, mentre un’altra no) e i protocolli appena visti tengono conto solo in parte di questi fattori. Per cui un’idea è quella di fornire un modo aggiuntivo per aggregare gerarchicamente le informazioni di routing in una rete Internet di grandi dimensioni. In questo modo si migliora la scalabilità. Si divide il problema dell’instradamento in due parti:
\begin{itemize}
  \item instradamento all’interno (within) di un singolo sistema autonomo
  \item instradamento tra (between) sistemi autonomi
\end{itemize}
Nel primo caso i router sono peer e hanno stesse regole e protocolli, nel secondo si devono creare delle gerarchie e domini, da cui si ottiene anche un altro nome per i sistemi autonomi (AS) in Internet: “domini di routing” (routing domains). La gerarchia di propagazione dei percorsi a due livelli segue due protocolli:
\begin{itemize}
  \item Protocollo di instradamento interdominio (standard per tutta Internet)
  \item Protocollo di routing intradominio (ogni AS sceglie il proprio)
\end{itemize}
\subsection{EGP e BGP}
Originariamente (quindi nel 1990 circa) il sistema Internet aveva una struttura ad albero, in cui vi era come root una grande rete backbone a cui poi si collegavano reti minori regionali o nazionali e poi nei livelli inferiori vi erano altre zone locali. La situazione è simile a quella del GARR, in cui il medesimo fa da backbone e i PoP sono le reti inferiori locali a cui si collegano le singole aziende/università/enti di ricerca. Questo tipo di struttura, tuttavia è rigida e andava bene finché Internet era piccola (sistema “accademico”), in cui vi era un’unica dorsale e i AS sono collegati solo come genitori e figli e non come peer, uno dei motivi che non ha permesso che la topologia diventasse generale col passare degli anni. Inoltre, i protocolli utilizzati al suo interno erano protocolli di routing inter-dominio, come il \textit{Protocollo Exterior Gateway (EGP)}. \\

Per cui la soluzione è stata quella di creare tanti provider di tipo backbone (es.: governativo, universitario, commerciale, ecc.) connessi tra loro, ottenendo una sorta di struttura a grafo. Da ricordare che vi possono essere anche provider minori (consumer ISP) che sfruttano delle backbone per offrire un servizio ai propri host. Con questa soluzione ci si può collegare a più backbone o provider (maggiore affidabilità). Per quanto riguarda il protocollo utilizzato, in questo caso si ha il \textit{Border Gateway Protocol (BGP)}. Esso presuppone che Internet sia un insieme arbitrariamente interconnesso di AS. Infatti l’Internet di oggi consiste in un’interconnessione di più reti dorsali (di solito sono chiamate reti di fornitori di servizi e sono gestite da aziende private piuttosto che dal governo). I siti sono connessi l’uno all’ altro in modi arbitrario. Alcune grandi aziende si collegano direttamente a una o più reti dorsali, mentre altre si collegano a fornitori di servizi più piccoli e non a delle dorsali. Molti fornitori di servizi esistono principalmente per fornire servizi ai “consumatori” (individui con PC nelle loro case), e questi fornitori devono connettersi ai fornitori di backbone. Spesso molti provider si organizzano per interconnettersi tra loro in un unico “punto di peering”.Il protocollo BGP non ha lo scopo di ottimizzare i percorsi, piuttosto ha quello di trovare i percorsi (privi di loop). Infatti, il protocollo BGP-4 per i gateway di confine, è atto a svolgere questa funzione. Si parte dall’ipotesi che Internet sia un insieme arbitrariamente interconnesso di AS. Si definisce traffico locale (local traffic) quello che ha origine o termina su nodi all’interno di un AS e traffico di transito (trasit traffic) quello che attraversa un AS. A questo punto si può classificare gli AS in tre tipi:
\begin{itemize}
  \item \textbf{AS stub}: un AS che ha una sola connessione con un altro AS; un AS di questo tipo trasporta solo traffico locale (la small corporation nella figura sopra).
  \item \textbf{AS multi-homed}: un AS che ha connessioni con più di un altro AS e a cui, però, non interessa trasportare traffico di transito (la large corporation in alto nella figura): semplicemente sfrutta uno o più link per motivi propri (es.: per affidabilità, efficienza, ecc.). Quindi, queste AS non sono provider, ma magari aziende.
  \item \textbf{Transit AS}: un AS che ha connessioni con più di un altro AS ed è progettato per trasportare sia il traffico di transito che quello locale (backbone provider). Infatti, il traffico che gira (entra/esce) in questi AS non sempre riguarda il traffico di questi AS stessi, ma si serve di essi per passare da una parte a un’altra della rete.
\end{itemize}
Come appena detto, BGP non ottimizza i percorsi, ma li trova soltanto e fa attenzione che non contengano cicli, poiché a queste dimensioni della rete conta la raggiungibilità piuttosto che l’ottimizzazione, trovare un percorso che si avvicini a quello ottimale è considerato un grande risultato. Infatti, gli obiettivi di BGP sono sostanzialmente:
\begin{itemize}
  \item Scalabilità: un router della dorsale Internet deve essere in grado di inoltrare qualsiasi pacchetto destinato a qualsiasi punto della rete. Dunque, deve avere una tabella di routing che fornisca una corrispondenza per qualsiasi indirizzo IP valido;
  \item Natura autonoma dei domini: è impossibile calcolare costi di percorso significativi per un percorso che attraversa più AS, soprattutto per il fatto che se si ha un costo di 1000 attraverso un provider potrebbe implicare un ottimo percorso per esso, ma per un altro potrebbe essere inaccettabile;
  \item Problemi di fiducia (trust): il fornitore A potrebbe non essere disposto a credere a certe pubblicità del fornitore B, ad esempio se si hanno due server di streaming e uno promette un delay più basso dell’altro si potrebbe essere più invogliati a sceglierlo piuttosto dell’altro, ma magari è solo una trovata di marketing e in realtà è più inaffidabile.
\end{itemize}
Ogni AS ha uno o più router speciali chiamati speaker BGP. Uno speaker BGP fa da portavoce ai suoi peer (altri router speaker BGP), ovvero comunica agli altri portavoce a esso collegati quali sono le reti locali interne (a quel AS), e eventualmente le reti non sue (AS di transito), e fornisce quindi informazioni sul percorso (es.: essi si dicono con chi possono comunicare in modo da “avere un’idea della rete” e sapere a chi
inoltrare i pacchetti per farli arrivare a destinazione). Oltre agli speaker BGP, un AS ha uno o più gateway di frontiera, che non devono necessariamente coincidere con gli speaker. I gateway di frontiera sono i router attraverso i quali i pacchetti entrano ed escono dal AS. Esistono due versioni del protocollo:
\begin{itemize}
  \item BGP esterno (eBGP): utilizzato tra router appartenenti a diversi AS
  \item BGP interno (iBGP): utilizzato tra router all’interno dello stesso AS
\end{itemize}
\subsection{Integrazione del routing interdominio e intradominio}
Per quanto riguarda il funzionamento vero e proprio si può immaginare un AS come una “nuvoletta”, in cui sono connessi più router, di cui alcuni di confine e comunicanti con l’esterno. Tutti i router eseguono iBGP e un protocollo di routing intra-dominio (es.: RIP, OSPF, ecc.). I router di confine (A, D, E) eseguono anche eBGP verso altri AS. Esiste una tabella di routing BGP valida per l’intero AS, le cui informazioni vengono poi propagate ai router interni tramite iBGP (in immagine).BGP funziona propagando i percorsi (path) completi e che portano a una certa destinazione, quindi hanno sempre delle reti definite.Dunque in un sistema completo si possono avere decine e decine di router, con decine di reti e decine di migliaia di host e non ce ne si accorge in quanto la rete sembra tutta omogenea. Inoltre, il BGP non appartiene a nessuna delle due classi principali di protocolli di routing (vettori di distanza e protocolli a stato di link). BGP pubblicizza percorsi completi come liste enumerate di AS per raggiungere una particolare rete. Le regole di inoltro e di accettazione dei pacchetti dipendono da come sono configurati gli speaker BGP, in quanto sta a loro sapere se accettare certi tipi di pacchetti (o con certi IP host) o meno.
\subsection{Prevenzione loop BGP e policy routing}
In realtà i router BGP non scartano i pacchetti solo in base alle proprie regole interne, ma anzi, quando essi ne ricevono uno di configurazione per i percorsi, analizzano tutto il path e decidono se scartare il pacchetto o meno per evitare loop. Per sapere se è in presenza di un loop controlla che non ci sia il proprio AS lungo il percorso comunicato dall’ altro speaker. Dunque, un router BGP ignora un percorso se: il suo AS compare in esso (per evitare i loop tra AS), oppure se viola qualche politica dell’ AS. Ad esempio, se si hanno quattro AS come in figura e un router R1 volesse comunicare che è in collegamento con la rete N1, lo fa attraverso R2, il quale sa di essere il “forwarder router” per la rete N1 poiché ne è messo in comunicazione. Per cui R2 comunica a R3 di essere in comunicazione con la rete N1 utilizzando il link che attraversa la propria rete interna AS2 e il router R1. A sua volta R3 comunica a R4 che è in comunicazione con N1 tramite rete interna, AS2 e AS1 (tramite rispettivi router). Tutto questo viene eseguito controllando prima le proprie policy e regole interne a ogni speaker BGP.
\subsection{Problemi BGP}
È importante che i numeri di AS trasportati in BGP devono essere univoci. Ad esempio, nell’immagine precedente, l’ AS 2 può riconoscersi nel percorso AS dell’esempio solo se nessun altro AS si identifica nello stesso modo. I numeri AS (ASN) sono numeri a 32 bit assegnati da un’autorità centrale (IANA e registri Internet regionali, RFC 4893) e sono circa 60.000 AS, al momento. Essi sono i nodi della rete che eseguono eBGP.
\section{IPv6}
È un protocollo di rete alternativo a IPv4. Tutto quello visto finora, continua a valere anche per IPv6, in quanto la grande differenza sta nel fatto che gli indirizzi IP non sono più da 4 byte, ma da 6 byte. Tra le ragioni che hanno portato a questa variante c’è in primis il fatto che gli indirizzi IPv4 sono terminati. Per cui è stato ridefinito il concetto di IP, utilizzando indirizzi a 128 bit, che forniscono uno spazio di indirizzamento dell’ ordine di $3*10^{38}$. Inoltre, dato che le cose da modificare sarebbero state tante, si è sfruttata l’ occasione per aggiungere in IPv6 servizi e funzionalità che IPv4 non aveva: multicast, servizio in tempo reale, autenticazione e sicurezza (un dispositivo IPv4 non era obbligato a fornire servizi di sicurezza),autoconfigurazione, frammentazione end-to-end, funzionalità di routing migliorate (compreso il supporto per gli host mobili, ovvero staccare da una rete un dispositivo, spostarsi e riagganciarlo altrove, mantenendo lo stesso IPv6 di prima senza dover riconfigurare), estensibilità.In IPv6 l’indirizzamento/routing è senza classi (simile a CIDR) e si usa la notazione: x:x:x:x:x:x:x:x (x = numero esadecimale a 16 bit). Gli 0 contigui sono compressi (es.: 47CD::A456:0124 oppure :::::1 per indicare sé stessi). L’assegnazione degli indirizzi è garantita dagli ISP oppure dagli RIR(Regional Internet Registries).
\subsection{Formato del pacchetto}
Per quanto riguarda i pacchetti utilizzati in IPv6, l’intestazione “base” è di 40 byte e si hanno 64 Kbyte di carico utile (payload, con estensioni). Le intestazioni di estensione (ordine fisso, lunghezza per lo più fissa) hanno frammentazione, instradamento della sorgente, autenticazione e sicurezza e molti altri. In realtà l’header di IPv6 è più semplice di quello di IPv4, perché è stato pensato per essere modulare. I primi 4 bit sono per la versione (come per IPv4), mentre gli 8 bit successivi (trafficClass) servono per un meccanismo simile al ToS, il flowLabel serve a identificare i pacchetti che fanno parte dello stesso flusso, il payloadLen è per la lunghezza del pacchetto (max 64 Kbyte), il nextHeader è una sorta di puntatore che dice dopo di un certo numero di bit si ha un altro header (o se è 0 non c’è) e il hopLimit è il TTL. Un’informazione che in IPv4 c’era, ma nell’intestazione base di IPv6 no è ad esempio il campo di frammentazione, che viene aggiunto nell’header aggiuntivo. In quest ultimo si trovano anche le opzioni sulla sicurezza o sul routing prefissato (che strada deve prendere il pacchetto, impostata manualmente).
\subsection{Priorità}
Il campo trafficClass definisce la priorità di un datagramma rispetto ad altri della stessa origine, ovvero è un suggerimento che si dà a un router per poter gestire meglio il pacchetto. È utile nel caso in cui alcuni datagrammi debbano essere abbandonati a causa di congestioni. I valori possibili sono:
\begin{itemize}
  \item 0 = Traffico generico (valore di default)
  \item 1 = Traffico di background (es.: NNTP), per traffici asincroni non prioritari (es.: backup o aggiornamento SW mentre si stanno eseguendo altri servizi in foreground come l’update di un SO)
  \item 2 = Traffico senza attesa (es.: SMTP), si usa per le e-mail in cui anche se un messaggio arriva 5 minuti più tardi non è troppo un problema (di solito)
  \item 3 e 5 = riservato
  \item 4 = Traffico con attesa (es.: FTP, HTTP), quando si sta aspettando la risposta di un servizio in tempi ragionevoli
  \item 6 = Traffico interattivo (es.: TELNET, SSH)
  \item 7 = Traffico di controllo (es.: OSPF, RIP, SNMP, ecc.), serve per controllare ad esempio lo stato degli switch (es.: per gestire congestioni)
  \item 8-15 = Traffico non dipendente dalla congestione (non deve essere interrotto), ordinato per ridondanza (es.: audio/video) e sono realtime (es.: dati di sensori), quindi con alta priorità
\end{itemize}
\subsection{Flow Label}
È un numero casuale a 20 bit, assegnato dalla sorgente. I datagrammi appartenenti allo stesso flusso sono etichettati con lo stesso valore. Questa etichetta può essere utilizzata dai router per implementare un servizio coerente a tutti i datagrammi dello stesso flusso (es.: stesso instradamento, oppure stessa qualità del servizio per il soft real-time, quindi l’utilizzo delle risorse riservate, che precedentemente erano riservate per mezzo di altri protocolli come RSVP, ecc.). Questo servizio è molto simile ai circuiti virtuali ed è compito dei router capire come gestire al meglio i pacchetti in base al servizio richiesto. Di solito si usa all’interno del AS e potrebbe implicare la riconfigurazione dei router appartenenti allo stesso AS.
\subsection{Transizione da IPv4 a IPv6}
IPv4 e IPv6 non sono compatibili. In realtà IPv4 dovrebbe essere sostituito da IPv6, anche se questo è molto difficile: troppi host da aggiornare, non può essere fatto in una volta sola, non può essere forzato “per legge”, ecc. Per cui esistono tre tecniche per supportare la transizione:
\subsubsection{Doppio stack di protocollo (dual protocol stack)}
Un host (o un router) può implementare entrambi i protocolli (IPv4 e IPv6). Questo è utilizzato nella maggior parte dei sistemi operativi moderni grazie alla sua versatilità (dato che IPv4 e IPv6 non sono compatibili). La decisione su quale dei due utilizzare viene presa quando si deve stabilire una connessione. Inoltre, si usano due protocolli a livello 3 in quanto l’host appartiene a due reti internet separate. Le applicazioni devono gestire esplicitamente entrambi gli indirizzi, il che comporta più lavoro per gli sviluppatori e i distributori (es.: vedere Apache httpd.conf).
\subsubsection{Tunneling}
Se due host usano entrambi IPv6 e devono attraversare un certo numero di router IPv4, è possibile creare “tunnel” o “ponte” da una regione IPv6 a un’altra, incapsulando i datagrammi IPv6 in datagrammi IPv4. Questo incapsulamento è dato dal fatto che gli indirizzi IPv6 sono pochi e isolati in mezzo a una marea di IPv4 e se un router IPv4 vede un pacchetto IPv6, non lo riconosce e lo scarta. Dunque il payload è un intero pacchetto IPv6 che viene fatto passare per il tunnel e il destinatario in questo caso diventa il router di destinazione (della rete dell’host) e non più l’host stesso. Questo approccio è gestito in modo trasparente dai router e gli host possono comunicare in IPv6, a discapito di perdere alcuni vantaggi di IPv6 (vengono persi nella regione IPv4), in più si crea un maggiore overhead e frammentazione (i pacchetti IPv6 sono più grandi e vengono incapsulati ulteriormente).
\subsubsection{Traduzione dell’intestazione}
Permette di trasmettere un datagramma IPv6 a un host IPv4. Il router di destinazione sostituisce l’header IPv6 con l’equivalente IPv4 e viceversa. Simile al NAT, ma la traduzione avviene tra indirizzi e intestazioni IPv4 e IPv6. Per fare ciò si usa un bridge (livello 3).
\section{Internet Multicast}
Questa funzionalità era presente anche in IPv4, solo che essendo un po’ difficile da implementare non viene quasi mai utilizzata. Come già spiegato, con multicast si intende la possibilità di inviare un messaggio a un gruppo di host riceventi, senza conoscerli o specificarli. Invece, con la sigla uno-a-molti (one-to-many) ci si riferisce a un multicast specifico della sorgente (SSM), quindi un host ricevente specifica un gruppo multicast in cui viene host specifico è colui che invia. Mentre con molti-a-molti (many-to-many) si intende un multicast di qualsiasi sorgente (ASM), ad esempio le teleconferenze multimediali, i giochi online multigiocatore, le simulazioni distribuite, ecc.
\subsection{Multicast IP}
Utilizzando il multicast IP, è possibile inviare lo stesso pacchetto a ogni membro del gruppo. Dunque un host invia una singola copia del pacchetto, indirizzato all’indirizzo multicast del gruppo, e poi di esso verrà data una copia per ogni singolo host membro del gruppo. L’host sorgente non ha bisogno di conoscere l’indirizzo IP unicast di ciascun membro, in quanto i pacchetti vengono duplicati dai router lungo il percorso, quando necessario.Se la rete fosse di grandi dimensioni, si dovrebbe configurare tanti router in modo che sappiano dove siano i vari indirizzi multicast,il che diventa impossibile se la rete fosse a livello globale. Tuttavia se non si avesse il supporto del multicast, una sorgente dovrebbe inviare un pacchetto separato con gli stessi dati a ciascun membro del gruppo e questa ridondanza consumerebbe più larghezza di banda. Inoltre il traffico ridondante non sarebbe distribuito in modo uniforme, ma si concentrerebbe vicino all’host mittente. In più, la sorgente dovrebbe tenere traccia dell’indirizzo IP di ogni membro del gruppo (potrebbero essere tanti!), quindi anche utilizzare solo unicast sarebbe ancora di più uno svantaggio. Per fortuna, IP fornisce un modello multicast molti-a-molti (a livello IP) basato su gruppi multicast. Ogni gruppo ha il proprio indirizzo IP multicast nella classe D (da 224.0.0.0 a 239.255.255.255). Gli indirizzi dei gruppi sono assegnati in diversi modi:
\begin{itemize}
  \item 224.0.0.0/24: l’uso è definito dalla IANA, ma è limitato alla rete locale.;
  \item 224.0.1.0/24: gruppi globali, assegnati staticamente dalla IANA (es.: 224.0.1.1 è NTP)
  \item Affittando dinamicamente per il tempo di una sessione, utilizzando il protocollo SAP/SDP (es.:per sessioni di videogiochi online multiplayer)
\end{itemize}
Per quanto riguarda la gestione dei gruppi multicast, un host può entrare e uscire dai gruppi come anche un host può far parte di più gruppi. Per cui un host segnala il suo desiderio di entrare o uscire da un gruppo multicast comunicando con il suo router locale e utilizzando uno speciale protocollo. In IPv4 si usava il protocollo di gestione dei gruppi Internet (IGMP), mentre in IPv6 c’è il Multicast Listener Discovery (MLD). Il router ha la responsabilità di far sì che il multicast si comporti correttamente nei confronti dell’host.
\begin{example}
  Ipotizziamo di avere una sorgente che sta già inviando un flusso di pacchetti ad un indirizzo multicast, un host effettu la richiesta di unione a tale gruppo multicast quindi invia un pacchetto IGMP al router locale. Il router inoltra la richiesta con l'indirizzo IP del multicast agli altri router vicini fino ad arrivare a un router intermeti che fa da streaming di tale sorgente. A quel punto i pacchetti della sorgente vengono inoltrati anche al nuovo host tramite i router intermedi el'host stesso può a sua volta inviare pacchetti al multicast a cui si è appena unito(tutto all'insaputa della sorgente). Se ad un certo punto l'host non volesse più ricevere pacchetti dal multicast, invia una richiesta ICMP al proprio router locale e questo la fa arrivare al router intermedio che fa da sorgente. Quest'ultimo capisce che non serve più inviare il flusso in quella direzione e smette di inviare su quel rmao della rete.
\end{example}

Altro aspetto da citare è la portata dei messaggi che può essere definita scegliento un TTL appropriato. Per quanto riguarda invece le questioni di sicurezza come la segretezza, integrità e autenticazione, ecc, non sono gestite a questo livello.

Le tabelle di inoltro unicast di un router indicano per qualsiasi indirizzo ip, quale link utilizzare per inoltrare il pacchetto unicast. Per supportare il multicast, un router deve anche disporre di tabelle di inoltro multicast che indicano, in base all' indirizzo di destinazione multicast, quali collegamenti utilizzare per inlotrare il pacchetto multicast. Tali tabelle specificano collettivamente una serie di alberi che vengono chiamati alberi di distribuzione multicast. L' instradamento multicast  è il processo con cui vengono determinati i multicast distribution tree.Data l'assenza di di criteri ben precisi su come costruire e popolare le tabelle in modo efficiente e scalabile, nel corso del tempo sono stati proposti molti protocolli di routing e due grandi famiglie di alberi: quelli basati sulla sorgente e quelli condivisi dal gruppo.
\subsection{Routing multicast basato sulla sorgente}
Utilizzando questo approccio ogni router deve avere un albero del percoso più breve per ogni gruppo. Ogni router quindi si costruisce un propria visione della rete tramite una tabella di inoltro (non contiene ercorsi completi) che presenta solo i router necessari al forward. 
\subsection{Routing multicast ad albero condiviso per gruppi}
Nell'approccio group-shared tree solo un router ha un albero del percorso più breve per ogni gruppo ed è coinvolto nel multicasting. Esso tiene traccia di tutto il traffico di pacchetti togliendo il compito agli altri (sarebbe ridondante).Per effettuare una richiesta di inserimento nel gruppo o per mandare un messaggio in multicast si deve inviare un pacchetto in unicast al router core che poi gestirà la richiesta. Rispetto alla soluzione precedente questa risulta più facile, con il difetto che non sempre è scalabile.
\subsection{Protocollo di instradamento multicast a vettore di distanza}\label{sec:Protocollo di instradamento multicast a vettore di distanza}
Il primo protocollo di routing multicast implementato è il \textit{Distance-Vector Multicast Routing Protocol}. Estende i protocolli di routing distance vector per il routing unicast e segue la strategia \textit{Flood and Prune} ovvero: inonda la rete di traffico multicast e nel frattempo pota i rami non interessati al traffico.
\subsubsection{Flooding:}
Per quanto riguarda il flooding vero e proprio, si cerca di inviare i messaggi in quantità giusta e lo si fa utilizzando una tecnica chiamata \textit{Reverse Path Flooding}.Lo scopo è quello di inviare, dalla sorgente, un certo quantitativo di pacchetti in modo tale da arrivare a tutti i nodi una sola volta.\textbf{NB Questo è il motivo per cui non si usa il Flooding puro in quando invierebbe a tutti molti pacchetti.}

Utilizzando gli algoritmi preesistenti (già visti in RIP per unicast) si sa quali sono i next hop da fare per arrivare a tutti i nodi della rete. Ogni router sà se deve inoltrare il pacchetto o no in base a quale indirizzo e porta sono stati utilizzati per fargli arrivare il pacchetto.A questo punto il router guarda il percorso e se vede che il percorso più breve è lo stesso di quello utilizzato per inviargli il pacchetto allora procede con l'inoltro su tutte le altre porte, altrimenti non forwarda nulla. Così facendo non ci sono loop durante il flooding. Tuttavia le reti grandi, quindi con un numero maggiore di router posso ricevere dei pacchetti duplicati, come soluzione viene assegnato ad un router,per ogni LAN, il ruolo di genitore; solitamente il genitore è il router che ha lo shortest-path verso la sorgente. I genitori sono gli unici che possono inoltrare i pacchetti.Si rompono poi i legami andando a scegliere il router con il minore ttl oppure con l'indirizzo ip più piccolo.Questa strategia si chiama \textit{Reverse Path Broadcast} e garantisce che ogni LAN riceva esattamente una copia di ogni pacchetto.
\subsubsection{Prune:}
In questa fase si potano le reti che non hanno host interessati al gruppo multicast, facendo così si limita ulteriormente il processo di flooding e non lo si espande a tutti gli host della rete. Il processo di pruning avvente in due fasi:
\begin{enumerate}
  \item Si determina se la LAN è una foglia dell'albero di distribuzione con nessun membro nel multicast. Una LAN è una foglia se il suo genitore è l'unico router della LAN.Tutti gli host che vogliono partecipare al multicast devono notificare al router utilizzando il protocollo IGMP, in questo modo il router sa se la LAN ha dei membri del multicast.
  \item Si propaga l'informazione "Nessun membro del multicast è qui",si aggiorna ogni voce <Destinazione,Costo>, nei vettori di distanza inviati ai vicini, con l'elenco dei gruppi per i quali questa rete è interessata a ricevere pacchetti multicast. In questo modo ogni vicino sà se deve considerare quel preciso router nel percorso inverso di flooding/broadcast. Tuttavia includere sempre l'elenco dei gruppi interessati nei vettori può risultare pesante e inutile (per esempio nel caso in cui una LAN sia interessata a diversi gruppi ma nessun host sta trasmettendo in quest'ultimi).
\end{enumerate}
Riassumento, quando inizia una trasmissione multicast sul gruppo:
\begin{itemize}
  \item all’inizio si inonda tutta la rete, tramite RPF/RPB, costruendo un albero di distribuzione che copre praticamente tutta la rete (perché nessuno dice ancora di non essere interessato a G);
  \item l’ albero poi viene potato ovvero i router non interessati al multicast iniziano a comunicare di non trasmettere più a loro;
\end{itemize}
Se un host di una rete potata/tagliata si unisce al gruppo multicast in un secondo momento, notifica il suo router locale utilizzando IGMP e il router inizia ad accettare il traffico dai suoi vicini (innesto = grafting). Nel complesso, il DVMRP funziona bene su piccola scala (intra-dominio) e meno su grande scala a causa di questo continuo flooding e pruning.
\subsection{Multicast indipendente dal protocollo}
Questo protocollo nasce in seguito al precedente ed è indipendente dalla costruzione della topologia di rete, infatti è basato su RIP.Il \textit{Protocol Indipendent Multicast} ha due modalità diverse:
\begin{itemize}
  \item PIM Dense Mode:Molto simile a \ref{sec:Protocollo di instradamento multicast a vettore di distanza} ma indipendente dal protocollo sottostante di routing unicast. E' basato sull'utilizzo degli alberi: ogni router possiede una copia del Multicast distribution Tree. Utilizza la strategia di Flooding e Pruning, in particolare RPF per la distribuzione del traffico e Flooding e Pruninig per gli alberi. Funziona molto bene quando la rete non è molto grande e ci sono pochi router che vengono attraversati.
  \item PIM Sparse Mode:può essere utilizzata per grandi gruppi di router, anche sparsi per internet. Per ogni gruppo viene costruito un \textit{Group-shared tree} e per ogni gruppo viene nominato un router che raggruppa gli altri (Rendez-vous router). Gli altri router possono entrare a far parte del grouppo inviando un richiesta di join al rendez-vous point. A mano a mano che la richiesta di join attraversa i vari router per arrivare al rendez-vous router crea il multicast distribution tree con radice il rendez-vous point. Ogni router analizza il messaggio e aggiunge alla sua tabella la regola per inoltrare e scaricare il il percorso fatto la join proveninte la gruppo multicast.Se il router non fa ancora parte del gruppo,
        \begin{enumerate}
          \item inoltra il join al RP (Router Principale) e segna l'interfaccia corrispondente come l'unica da cui può provenire il traffico
          \item non fa nulla
        \end{enumerate}
    In questa maniera i router “imparano” la strada da far seguire ai pacchetti multicast per fare in modo che arrivino a destinazione. Essi segnano nelle proprie tabelle di inoltro il percorso, fino a quando non scade la richiesta dell’host di restare nel gruppo multicast. Così quando si aggrega un altro host al flusso multicast, i router hanno già la strada segnata.
\end{itemize}
\subsection{Backbone Multicast (MBONE)}
Nella realtà solo una piccola parte dei router implementa il routing multicast a causa del traffico molto costoso e molto spesso questi router non sono contigui, cioè sono collegati con regione che non sono multicast-aware. Una soluzione temporanea può essere quella di creare dei tunnel unicast tra router multicast, attraversando router non multicast.I tunnel sono collegamenti logici e agiscono come una dorsale per il multicast, questa dorsale è chiamata \textbf{MBONE}. I pacchetti multicast sono incapsulati all'interno di normali pacchetti IP unicast e possono passare attraverso i router multicast-unaware.

C'è però da dire che anche MBONE non è molto diffiso come soluzione dato che ha diversi problemi di accounting, si aggiunge carico sul router e molto altro ancora.Come se non bastasse MBONE è obsoleto e sta per essere sostituito con IPv6 e PIMv6. Il multicast viene adottato all'interno di singole organizzazioni come per esempio dentro un Hotel o in un ospedale.
